\renewcommand{\Anno}{1998}

\renewcommand{\NumeroProblema}{3}

\renewcommand{\Testo}{I vertici di un triangolo nel piano cartesiano
hanno tutti coordiante razionali.  Dimostrare che l'area del triangolo
è un numero razionale.}

\renewcommand{\Risposta}{
% SSSUP-BOZZA-AI: Risposta; autore=ChatGPT; data=2026-11-08; stato=DA-VERIFICARE
% SSSUP-PROVENIENZA: AI: elaborazione nuova dal testo
Se i vertici sono $(x_1,y_1)$, $(x_2,y_2)$ e $(x_3,y_3)$, l'area è
\[
\frac12\left|(x_2-x_1)(y_3-y_1)-(x_3-x_1)(y_2-y_1)\right|\in\mathbb Q.
\]
}

\renewcommand{\Aiutino}{Vi sono diversi approcci possibili.
\begin{itemize}
	\item  Iniziare a considerare il caso in cui si tratta di un triangolo rettangolo con i cateti paralleli agli assi, poi ridurre il caso generale a somma e/o differenza di triangoli rettangoli.

	\item  Usare la geometria analitica: scelto un lato come base, basta calcolare la sua lunghezza e quella dell'altezza relativa.

	\item  Ridursi mediante un'omotetia al caso in cui le coordinate sono intere, quindi applicare il teorema di Pick.
\end{itemize}
}

\renewcommand{\SoluzioneUfficiale}{
}

\renewcommand{\SoluzioniAlternative}{
}

\renewcommand{\Osservazioni}{
}

\renewcommand{\Commenti}{
}
% Aggiunta redazionale: la soluzione preesistente non viene modificata.
\providecommand{\SoluzioneRedazionale}{}
\renewcommand{\SoluzioneRedazionale}{
% SSSUP-BOZZA-AI: SoluzioneRedazionale; autore=ChatGPT; data=2026-11-08; stato=DA-VERIFICARE
% SSSUP-PROVENIENZA: AI: elaborazione nuova dal testo
Indichiamo i vertici con $A=(x_1,y_1)$, $B=(x_2,y_2)$ e $C=(x_3,y_3)$, dove tutte le coordinate sono razionali. Il parallelogrammo generato dai vettori $B-A$ e $C-A$ ha area uguale al valore assoluto del loro determinante; il triangolo ne occupa la metà. Quindi
\[
\operatorname{Area}(ABC)=\frac12\left|(x_2-x_1)(y_3-y_1)-(x_3-x_1)(y_2-y_1)\right|.
\]
Differenze, prodotti, somma, valore assoluto e divisione per $2$ conservano la razionalità. L'area è pertanto razionale. La formula comprende anche il caso degenere, nel quale l'area è zero.
}
